\documentclass[tikz,border=10pt]{standalone}
\usepackage[no-math]{fontspec}
\usepackage[brazil]{babel}
\usepackage{amsmath,amssymb}
\IfFileExists{C:/Windows/Fonts/Lato_Regular.ttf}{%
  \setmainfont{Lato}[
    Path = C:/Windows/Fonts/,
    UprightFont = Lato_Regular.ttf,
    BoldFont = Lato_Bold.ttf,
    ItalicFont = Lato_Italic.ttf,
    BoldItalicFont = Lato_BoldItalic.ttf
  ]
  \setsansfont{Lato}[
    Path = C:/Windows/Fonts/,
    UprightFont = Lato_Regular.ttf,
    BoldFont = Lato_Bold.ttf,
    ItalicFont = Lato_Italic.ttf,
    BoldItalicFont = Lato_BoldItalic.ttf
  ]
}{%
  \RequirePackage[default]{lato}
}
\renewcommand{\familydefault}{\sfdefault}

\usepackage{xcolor}
\usepackage{tikz}
\usepackage{pgfplots}
\pgfplotsset{compat=1.18}
\usetikzlibrary{arrows.meta,positioning,shapes.geometric,fit,backgrounds,calc,decorations.pathreplacing,patterns}

% Paleta institucional UEFS / Bioengenharia de Solos
\definecolor{bsNavy}{HTML}{1F4268}
\definecolor{bsNavyDark}{HTML}{152E4A}
\definecolor{bsGold}{HTML}{B38A3D}
\definecolor{bsGoldDark}{HTML}{8C6927}
\definecolor{bsTeal}{HTML}{0E7C7B}
\definecolor{bsForest}{HTML}{2D6A4F}
\definecolor{bsAlert}{HTML}{9E2A2B}
\definecolor{bsEarth}{HTML}{8C5A3C}
\definecolor{bsWater}{HTML}{2E86AB}
\definecolor{bsInk}{HTML}{26313D}
\definecolor{bsMuted}{HTML}{5B6573}

\definecolor{bgNavy}{HTML}{EDF2F7}
\definecolor{bgGold}{HTML}{F8F3E9}
\definecolor{bgTeal}{HTML}{EBF6F5}
\definecolor{bgForest}{HTML}{EEF5F0}
\definecolor{bgAlert}{HTML}{FDF2F2}

\tikzset{
  >=Stealth,
  bsBox/.style={
    draw=#1, fill=white, rounded corners=5pt, line width=1.15pt,
    inner sep=7pt, align=center, font=\footnotesize\sffamily, text=bsInk
  },
  bsTint/.style={
    draw=#1, fill=#1!8!white, rounded corners=5pt, line width=1.15pt,
    inner sep=7pt, align=center, font=\footnotesize\sffamily, text=bsInk
  },
  bsBadge/.style={
    fill=#1, text=white, rounded corners=9pt, inner xsep=8pt, inner ysep=3pt,
    font=\scriptsize\bfseries\sffamily
  },
  bsTitle/.style={font=\normalsize\bfseries\sffamily, text=bsNavy, align=center},
  bsSub/.style={font=\footnotesize\sffamily, text=bsMuted, align=center},
  bsArrow/.style={->, >=Stealth, line width=1.35pt},
  bsDArrow/.style={<->, >=Stealth, line width=1.35pt}
}


\begin{document}
\begin{tikzpicture}[font=\sffamily]

% ============================================================
% TÍTULO GERAL
% ============================================================
\node[bsTitle, anchor=west] at (0, 7.80)
  {Transferência Gradual de Resistência no Ciclo de Vida do Gabião Vivo};
\node[bsSub, anchor=west] at (0, 7.35)
  {Substituição progressiva da contenção metálica pela matriz radicular autossustentável};

% ============================================================
% GRÁFICO PGFPLOTS
% ============================================================
\begin{axis}[
  at={(0.2cm,0.2cm)},
  width=15.0cm,
  height=6.7cm,
  xmin=0, xmax=15,
  ymin=0, ymax=108,
  xtick={0,3,6,9,12,15},
  xticklabels={0, 3, 6, 9, 12, 15+},
  ytick={0,25,50,75,100},
  yticklabels={0\%, 25\%, 50\%, 75\%, 100\%},
  xlabel={Tempo após implantação (Anos)},
  ylabel={Contribuição à Estabilidade (\%)},
  xlabel style={font=\footnotesize\bfseries\sffamily, text=bsInk},
  ylabel style={font=\footnotesize\bfseries\sffamily, text=bsInk},
  tick label style={font=\scriptsize\sffamily, text=bsMuted},
  axis lines=left,
  axis line style={bsInk, line width=0.9pt},
  ymajorgrids=true,
  grid style={dotted, bsMuted!45},
  clip=false,
  legend style={
    at={(0.98, 0.36)},
    anchor=south east,
    draw=bsNavy!30,
    fill=white,
    rounded corners=4pt,
    font=\scriptsize\sffamily,
    inner sep=5pt,
    cells={anchor=west}
  }
]

  % Faixas verticais das 3 fases
  \fill[bgNavy] (axis cs:0,0) rectangle (axis cs:3,100);
  \fill[bgGold] (axis cs:3,0) rectangle (axis cs:9,100);
  \fill[bgForest] (axis cs:9,0) rectangle (axis cs:15,100);

  \draw[bsMuted!60, dashed, line width=0.8pt] (axis cs:3,0) -- (axis cs:3,102);
  \draw[bsMuted!60, dashed, line width=0.8pt] (axis cs:9,0) -- (axis cs:9,102);

  % Marcadores de fase no topo (com cantos limpos e tamanho ajustado às faixas)
  \node[bsBadge=bsNavy, rounded corners=5pt, font=\tiny\bfseries\sffamily, inner xsep=6pt, anchor=south]
    at (axis cs:1.5, 102) {Fase 1 · Contenção Mecânica};
  \node[bsBadge=bsGold, rounded corners=5pt, font=\tiny\bfseries\sffamily, inner xsep=7pt, anchor=south]
    at (axis cs:6.0, 102) {Fase 2 · Transição Híbrida};
  \node[bsBadge=bsForest, rounded corners=5pt, font=\tiny\bfseries\sffamily, inner xsep=7pt, anchor=south]
    at (axis cs:12.0, 102) {Fase 3 · Estabilidade Viva};

  % Curva 1: Resistência Estrutural da Gaiola / Aço (tracejada)
  \addplot[
    color=bsNavy,
    dashed,
    line width=1.8pt,
    smooth
  ] coordinates {
    (0, 88) (1.5, 87) (3, 84) (5, 72) (6.2, 60) (7.5, 48) (9, 38) (12, 28) (15, 24)
  };
  \addlegendentry{Resistência Estrutural da Gaiola / Aço}

  % Curva 2: Resistência Biológica da Matriz Solo-Raiz (sólida verde)
  \addplot[
    color=bsForest,
    line width=2.2pt,
    smooth
  ] coordinates {
    (0, 4) (1.5, 12) (3, 28) (5, 50) (6.2, 60) (7.5, 73) (9, 84) (12, 92) (15, 95)
  };
  \addlegendentry{Resistência Biológica da Matriz Solo–Raiz}

  % Ponto de interseção e callout em área livre abaixo da interseção
  \filldraw[fill=bsGold, draw=bsNavyDark, line width=1.1pt] (axis cs:6.2, 60) circle (3.5pt);
  \node[
    draw=bsGoldDark, fill=white, rounded corners=4pt, line width=1.0pt,
    font=\scriptsize\bfseries\sffamily, text=bsNavyDark, align=center,
    inner sep=5pt, anchor=north
  ] (callout) at (axis cs:6.0, 26) {Ponto de Autossustentabilidade\\Biológica};
  \draw[bsArrow, bsGoldDark, line width=1.1pt] (callout.north) -- (axis cs:6.18, 57.5);

\end{axis}

\end{tikzpicture}
\end{document}
